\scanpage{26}
\mspage{\handwritten{17}}
conclusions de 7.2, sous la forme \unsure{affaablie|affaiblie} qu'on peut \deleted{résoudre} satisfaire à la congruence envisagée pour au moins deux valeurs de $a$ distinctes mod. $p$, et par suite pour avoir la conclusion de 7.3 (qui n'utilise essentiellement que 7.2 sous cette forme affaiblie).

\rlap{\begin{tikzpicture}[baseline=0pt,x=1cm,y=1cm]\draw[line width=.4pt] (-.72,-.15)--(-.08,-.15) .. controls (-.14,.13) and (-.14,.48) .. (-.08,.63)--(.12,.63);\draw[line width=.3pt] (.15,.63)--(.32,.62);\draw[line width=.22pt,opacity=.55] (.35,.62)--(.62,.61);\draw[line width=.3pt] (.74,.62)--(.78,.61);\draw[line width=.4pt] (.91,.62)--(1.02,.63);\end{tikzpicture}}\inserted{b)} Lorsqu'on veut 7.6 avec un $N$ donné, on \deleted{peut aussi} pourra sans doute trouver souvent une meilleure majoration de $d'_o$ que $(d_1\handwritten{+}1)(n+1)$, en décomposant $N$ en $N'+N''$, en choisissant $d'_1$ et $d''_1$ de façon que $N'\handwritten{\geq}rd_1'{}^n$, $N''\handwritten{\geq}rd_2'{}^n$, et construisant une forme produit de degré \rlap{\begin{tikzpicture}[baseline=0pt,x=1cm,y=1cm]\draw[line width=.4pt] (0,-.2) .. controls (-.05,.05) and (.25,.25) .. (.48,.29) .. controls (.82,.33) and (1.13,.45) .. (1.1,.59)--(1.04,.69);\end{tikzpicture}}\inserted{$d=$} $d'_1\unsure{d''_2|d'_2}+d''_1d''_2$ ($d'_{\unsure{\deleted{1}2|\deleted{2}1}},d''_2$ des entiers arbitraires $\handwritten{\geq}n+1$) dont l'ensemble des zéros communs avec $X$ rationnels\rlap{\begin{tikzpicture}[baseline=0pt,x=1cm,y=1cm]\draw[line width=.4pt] (.04,.38) .. controls (.17,.38) and (.28,.4) .. (.29,.5)--(.3,.76);\end{tikzpicture}} sur $k$ soit la somme disjointe de deux ensembles de cardinaux respectifs $N'$ et $N''$. Quitte à choisir $d'_1$ et $d''_1$ \deleted{assez gra} premiers entre eux, on trouve \rlap{\begin{tikzpicture}[baseline=0pt,x=1cm,y=1cm]\draw[line width=.4pt] (0,-.22) .. controls (-.1,-.05) and (.02,.22) .. (.3,.4) .. controls (.58,.59) and (.74,.7) .. (.74,.82);\end{tikzpicture}}\inserted{que} tout entier assez grand $d$ \rlap{\begin{tikzpicture}[baseline=0pt,x=1cm,y=1cm]\draw[line width=.4pt] (0,-.24) .. controls (-.12,-.1) and (.04,.17) .. (.36,.23) .. controls (.77,.32) and (1.47,.24) .. (1.72,.53) .. controls (1.8,.62) and (1.78,.66) .. (1.74,.7);\end{tikzpicture}}\inserted{se met} sous la forme envisagée.

c) On fera attention que les formes auxiliaires sans zéro non trivial utilisées dans la démonstration de 7.6 sont construites avec des formes normiques, et sont toujours décomposables géométriquement en facteurs linéaires. Ceci précise la remarque faite au début de la démonstration de 7.1 concernant la nécessité de travailler avec le gros ouvert $U_d$ tout entier.

Indiquons pour terminer la raison générale pour 7.5 :

\uline{Lemme} 7.8. Soient $S$ un schéma\rlap{\begin{tikzpicture}[baseline=0pt,x=1cm,y=1cm]\draw[line width=.35pt] (0,-.26)--(.08,-.12);\draw[line width=.28pt] (.105,-.07)--(.16,.035);\draw[line width=.25pt] (.18,.08)--(.22,.16);\draw[line width=.3pt] (.235,.205)--(.29,.32);\draw[line width=.28pt] (.31,.36)--(.35,.45);\end{tikzpicture}}\margin{\rlap{\begin{tikzpicture}[baseline=0pt,x=1cm,y=1cm]\draw[line width=.45pt] (-.32,-.3) .. controls (-.16,.03) and (.12,.42) .. (.18,.69);\end{tikzpicture}}\handwritten{normal}} de car. $p>0$, $\underline{V}$ un Module localement libre de type fini sur $S$ muni d'un \rlap{\begin{tikzpicture}[baseline=0pt,x=1cm,y=1cm]\draw[line width=.4pt] (0,-.16) .. controls (-.03,.1) and (.05,.24) .. (.36,.32) .. controls (.8,.43) and (1.15,.42) .. (1.26,.74);\end{tikzpicture}}\inserted{\handwritten{endo}}morphisme $p$-linéairex $f$, $t$ un point de $s$ et $s$ une spécialisation de $t$. Considérons les représentations des groupes $G_t=\operatorname{Gal}(\overline{k(t)}/k(t))$ et $G_s=\operatorname{Gal}(\overline{k(s)}/k(s))$ dans des vectoriels de dimension finie $V_1$ resp. $V_o$ sur le corps premier $\underline{F}_p$, définis par les parti s semi-simples des fibres \deleted{de} réduites de $\underline{V}$ en $t$ resp. en $s$. Alors il existe un sous-espace $V'_o$ de $V_1$, stable par un sous-groupe ouvert $G'_t$ de $G_t$, tel \rlap{\begin{tikzpicture}[baseline=0pt,x=1cm,y=1cm]\draw[line width=.4pt] (0,-.2) .. controls (-.04,-.02) and (.18,.24) .. (.49,.46);\draw[line width=.25pt,opacity=.7] (.68,.53) .. controls (.76,.55) and (.84,.58) .. (.9,.63);\end{tikzpicture}}\inserted{que} la représentation de l'image $G'_{t\,o}$ de $G'_t$ dans \deleted{$\unsure{V'|V^+}$} $\operatorname{Aut}(V'_o)$ par automorphismes de $V'_o$ soit isomorphe à celle de l'image $\unsure{G_{s\,o}|G_s,}$
% archive sheet number: 25
